\subsection{Relation to the case of maps of base spaces as morphisms}

Section~\ref{chap:functoriality} and Section~\ref{chap:operator} put different
morphisms on the same objects and obtained a reflection theorem of the same
shape for each.
In this subsection we clarify the relation between the two.

We state the conclusion first.
The functor $\mathcal{E}$ commutes with the two reflections
(Proposition~\ref{prop:E-commutes-with-D}), and as a consequence
\textbf{the uniqueness part of the main theorem of
Section~\ref{chap:functoriality} follows from the theorem of
Section~\ref{chap:operator}}
(Corollary~\ref{cor:ch3-uniqueness-from-ch4}).
The existence part, however, \textbf{does not follow}.
That is, one cannot say from the operator version alone that the extension comes
from a map of base spaces
(Remark~\ref{rem:two-chapters-relation}).

\subsubsection{A functor between the ae-categories}

\begin{lemma}\label{lem:E-descends-to-ae}
  For a morphism of integration systems
  $F:(X,L,I)\to(Y,M,J)$ put
  \[
    \mathcal{E}_{\mathrm{ae}}(F):=\pi_J\circ F^*:L\to\widetilde{M} .
  \]
  Then the following hold.
  \begin{enumerate}
    \item $\mathcal{E}_{\mathrm{ae}}(F)$ is an operator morphism into the
      quotient.
    \item $F=_{\mathrm{ae}}G$ and
      $\mathcal{E}_{\mathrm{ae}}(F)=\mathcal{E}_{\mathrm{ae}}(G)$ are
      equivalent.
    \item The assignment $[F]\mapsto\mathcal{E}_{\mathrm{ae}}(F)$ defines a
      \textbf{faithful} functor
      \[
        \mathcal{E}_{\mathrm{ae}}:
        \mathbf{Int}_{\mathrm{ae}}\to\mathbf{Int}_{\mathrm{op},\mathrm{ae}} ,
      \]
      which is the identity on objects.
  \end{enumerate}
\end{lemma}
\begin{proof}
  (1) By Lemma~\ref{lem:base-map-gives-operator} the map $F^*$ is linear and
  positive and preserves the integral, and $\pi_J$ is a linear Riesz
  homomorphism (Lemma~\ref{lem:quotient-riesz}), so the composite is linear and
  positive as well.
  Moreover $\widetilde{J}(\pi_J(F^*f))=J(F^*f)=I(f)$.

  (2) $\pi_J(F^*f)=\pi_J(G^*f)$ is equivalent to
  $F^*f-G^*f\in\mathcal{N}_J$, that is, to
  $J(|F^*f-G^*f|)=0$.
  That this holds for every $f\in L$ was the definition of
  $F=_{\mathrm{ae}}G$ (Definition~\ref{def:ae-equal}).

  (3) The ``only if'' direction of (2) shows that the assignment is independent
  of the choice of representative.
  The converse direction gives injectivity, that is, faithfulness.
  For identity morphisms,
  $\mathcal{E}_{\mathrm{ae}}(\operatorname{id})=\pi_I$, which is the identity
  morphism of $\mathbf{Int}_{\mathrm{op},\mathrm{ae}}$
  (Definition~\ref{def:int-op-ae}).
  As for composition, let $F:(X,L,I)\to(Y,M,J)$ and
  $G:(Y,M,J)\to(Z,N,K)$. By the definition of composition in
  Definition~\ref{def:int-op-ae}, the map
  $\widetilde{M}\to\widetilde{N}$ induced by
  $\mathcal{E}_{\mathrm{ae}}(G)$ is $\pi_J(h)\mapsto\pi_K(G^*h)$. Hence, for
  every $f\in L$,
  \[
    \bigl(\widetilde{\mathcal{E}_{\mathrm{ae}}(G)}\circ\mathcal{E}_{\mathrm{ae}}(F)\bigr)(f)
    =\pi_K\bigl(G^*(F^*f)\bigr)
    =\pi_K\bigl((G\circ F)^*f\bigr)
    =\mathcal{E}_{\mathrm{ae}}(G\circ F)(f) .
  \]
\end{proof}

\begin{remark}\label{rem:E-ae-faithful-not-full}
  As stated at the beginning of Section~\ref{chap:operator}, the functor
  $\mathcal{E}:\mathbf{Int}\to\mathbf{Int}_{\mathrm{op}}$ before passing to the
  ae-quotient need not be faithful, because distinct maps of base spaces can
  give the same pullback.
  Once the ae-quotient is taken, exactly that difference is collapsed by
  $=_{\mathrm{ae}}$, so $\mathcal{E}_{\mathrm{ae}}$ becomes faithful.

  On the other hand $\mathcal{E}_{\mathrm{ae}}$ is \textbf{not full}.
  Only the morphisms coming from lattice homomorphisms, that is, the
  deterministic ones, lie in the image, and a general operator morphism is not
  among them.
  In the case of finite sets, the image consists only of those
  column-stochastic matrices each of whose rows has exactly one $1$.
\end{remark}

\subsubsection{Commutativity of the two reflections}

We write $\mathcal{D}_{\mathrm{op}}$ for the left adjoint of
$\iota_{B,\mathrm{op}}$ given by
Theorem~\ref{thm:operator-main-theorem}.
Its action is as stated in Theorem~\ref{thm:operator-main-theorem}.

\begin{proposition}\label{prop:E-commutes-with-D}
  The following diagram commutes.
  \[
    \begin{array}{ccc}
      \mathbf{Int}_{\mathrm{ae}}
        & \xrightarrow{\ \mathcal{D}_{\mathrm{ae}}\ }
        & \mathbf{DInt}_{\mathrm{ae}}\\[2pt]
      \big\downarrow{\scriptstyle\mathcal{E}_{\mathrm{ae}}}
        & & \big\downarrow{\scriptstyle\mathcal{E}_{\mathrm{ae}}}\\[2pt]
      \mathbf{Int}_{\mathrm{op},\mathrm{ae}}
        & \xrightarrow{\ \mathcal{D}_{\mathrm{op}}\ }
        & \mathbf{BInt}_{\mathrm{op},\mathrm{ae}}
    \end{array}
  \]
  That is,
  $\mathcal{E}_{\mathrm{ae}}\circ\mathcal{D}_{\mathrm{ae}}
  =\mathcal{D}_{\mathrm{op}}\circ\mathcal{E}_{\mathrm{ae}}$.
\end{proposition}
\begin{proof}
  On objects both are the Daniell extension
  $(X,L,I)\mapsto(X,L^1,I^1)$, and $\mathcal{E}_{\mathrm{ae}}$ does not change
  objects, so they agree.

  We argue for morphisms. Let $[F]:(X,L,I)\to(X',L',I')$ be a morphism of
  $\mathbf{Int}_{\mathrm{ae}}$.
  Write the injective Riesz homomorphism determined by
  Lemma~\ref{lem:quotient-embedding} as
  \[
    j':\widetilde{L'}\longrightarrow\widetilde{L'^1} .
  \]
  The unit morphism
  $i_{A'}:(X',L',I')\to(X',L'^1,I'^1)$ in
  $\mathbf{Int}_{\mathrm{op},\mathrm{ae}}$ is
  $\pi_{I'^1}|_{L'}=j'\circ\pi_{I'}$.

  Since $\mathcal{D}_{\mathrm{op}}$ is a left adjoint, for an operator morphism
  into the quotient $\Lambda:L\to\widetilde{L'}$ the morphism
  $\mathcal{D}_{\mathrm{op}}(\Lambda)$ is the unique operator morphism into the
  quotient $L^1\to\widetilde{L'^1}$ satisfying
  \[
    \mathcal{D}_{\mathrm{op}}(\Lambda)\bigr|_{L}=j'\circ\Lambda .
  \]
  Applying this to
  $\Lambda=\mathcal{E}_{\mathrm{ae}}[F]=\pi_{I'}\circ F^*$ gives the composite
  going down then across in the diagram, which we write as $\Phi_2$.

  The composite going across then down, on the other hand, is
  $\mathcal{E}_{\mathrm{ae}}$ applied to
  $\mathcal{D}_{\mathrm{ae}}[F]=[F^1]$, that is,
  \[
    \Phi_1:=\pi_{I'^1}\circ(F^1)^*:L^1\to\widetilde{L'^1} .
  \]
  Applying Lemma~\ref{lem:E-descends-to-ae} to the morphism $F^1$, the map
  $\Phi_1$ is an operator morphism into the quotient.

  We check that $\Phi_1$ satisfies the same characterization as $\Phi_2$.
  For $f\in L$ we have $(F^1)^*f=F^*f\in L'$, so the equality
  $\pi_{I'^1}|_{L'}=j'\circ\pi_{I'}$ of
  Lemma~\ref{lem:quotient-embedding} gives
  \[
    \Phi_1(f)
    =\pi_{I'^1}(F^*f)
    =j'\bigl(\pi_{I'}(F^*f)\bigr)
    =\bigl(j'\circ\mathcal{E}_{\mathrm{ae}}[F]\bigr)(f) .
  \]
  That is, $\Phi_1|_L=j'\circ\mathcal{E}_{\mathrm{ae}}[F]=\Phi_2|_L$.

  The system $(X',L'^1,I'^1)$ is a Daniell closed integration system
  (Lemma~\ref{lem:daniell-ext-is-complete}), and
  Lemma~\ref{lem:quotient-is-banach} shows that $\widetilde{L'^1}$ is a Banach
  space, so the \textbf{uniqueness} in
  Theorem~\ref{thm:operator-extension} applies and we obtain
  $\Phi_1=\Phi_2$.
\end{proof}

\begin{remark}
  Note that this proof relies on the uniqueness of
  Section~\ref{chap:operator}.
  Within the framework of Section~\ref{chap:functoriality} alone there is no
  means of saying that the two extensions agree, because they are morphisms of
  different categories.
\end{remark}

\subsubsection{Relation between the two main theorems}\label{sec:two-chapters}

\begin{corollary}\label{cor:ch3-uniqueness-from-ch4}
  Let $A$ be an object of $\mathbf{Int}_{\mathrm{ae}}$, let $B$ be an object of
  $\mathbf{DInt}_{\mathrm{ae}}$, and let
  $[H],[H']:\mathcal{D}_{\mathrm{ae}}A\to B$ be morphisms.
  Then
  \[
    \iota_{D,\mathrm{ae}}([H])\circ[i_A]
    =\iota_{D,\mathrm{ae}}([H'])\circ[i_A]
    \quad\Longrightarrow\quad
    [H]=[H']
  \]
  holds. That is, the uniqueness part of
  Theorem~\ref{thm:main-ae} follows from
  Theorem~\ref{thm:operator-extension} of
  Section~\ref{chap:operator}.
\end{corollary}
\begin{proof}
  Apply $\mathcal{E}_{\mathrm{ae}}$ to both sides.
  By Lemma~\ref{lem:E-descends-to-ae} the assignment
  $\mathcal{E}_{\mathrm{ae}}$ is a functor, and
  $\mathcal{E}_{\mathrm{ae}}[i_A]$ is the morphism determined by the inclusion
  $L\subset L^1$, so $\mathcal{E}_{\mathrm{ae}}[H]$ and
  $\mathcal{E}_{\mathrm{ae}}[H']$ are operator morphisms into the quotient
  $L^1\to\widetilde{M}$ agreeing on $L$.
  Since $B$ is Daniell closed, $\widetilde{M}$ is a Banach space by
  Lemma~\ref{lem:quotient-is-banach}, and the uniqueness in
  Theorem~\ref{thm:operator-extension} gives
  $\mathcal{E}_{\mathrm{ae}}[H]=\mathcal{E}_{\mathrm{ae}}[H']$.
  As $\mathcal{E}_{\mathrm{ae}}$ is faithful
  (Lemma~\ref{lem:E-descends-to-ae}), we obtain $[H]=[H']$.
\end{proof}

\begin{remark}[relation between the two sections]\label{rem:two-chapters-relation}
  By contrast, the \textbf{existence} part of
  Theorem~\ref{thm:main-ae} does not follow from the theorem of
  Section~\ref{chap:operator}.
  Given $[F]:A\to\iota_{D,\mathrm{ae}}B$,
  Theorem~\ref{thm:operator-extension} supplies an extension of
  $\mathcal{E}_{\mathrm{ae}}[F]$ as a morphism of
  $\mathbf{Int}_{\mathrm{op},\mathrm{ae}}$, but it does not guarantee that this
  lies in the image of $\mathcal{E}_{\mathrm{ae}}$, that is, that it comes from
  a map of base spaces, because $\mathcal{E}_{\mathrm{ae}}$ is not full
  (Remark~\ref{rem:E-ae-faithful-not-full}).
  To obtain existence one needs the construction specific to
  Section~\ref{chap:functoriality}
  (Lemma~\ref{lem:daniell-ext-is-functor-1}), by which the extension is
  obtained without changing the map of base spaces.

  The contributions of the two theorems divide as follows.

  \begin{center}
  \begin{tabular}{@{}lll@{}}
    \toprule
     & Existence & Uniqueness \\
    \midrule
    Section~\ref{chap:functoriality}, non-ae & needs its own construction & follows automatically (no content) \\
    Section~\ref{chap:functoriality}, ae & as above & by density \\
    Section~\ref{chap:operator} (operators) & not obtainable from the operator version & by contractivity and density \\
    \bottomrule
  \end{tabular}
  \end{center}

  Care is needed in how this is phrased.
  Since $\mathcal{E}_{\mathrm{ae}}$ is faithful,
  $\mathbf{Int}_{\mathrm{ae}}$ is embedded (without changing objects) in
  $\mathbf{Int}_{\mathrm{op},\mathrm{ae}}$.
  In that sense \textbf{the category of
  Section~\ref{chap:functoriality} is a wide subcategory of the category of
  Section~\ref{chap:operator}}, and it is not accurate to say that ``neither is
  a special case of the other''.
  Nor is Section~\ref{chap:functoriality} in the relation of completing a single
  theorem together with Section~\ref{chap:operator}: it proves both existence
  and uniqueness on its own.
  The relations that can be stated precisely are the following.

  \begin{itemize}
    \item $\mathcal{E}_{\mathrm{ae}}$ is faithful but not full
      (Remark~\ref{rem:E-ae-faithful-not-full}).
    \item The two reflections are compatible under this functor
      (Proposition~\ref{prop:E-commutes-with-D}).
    \item The \emph{uniqueness} of the base-space-map version can be derived
      from the operator version
      (Corollary~\ref{cor:ch3-uniqueness-from-ch4}).
    \item But \textbf{that the extension stays in the image of
      $\mathcal{E}_{\mathrm{ae}}$, that is, that it comes from a map of base
      spaces, cannot be derived from the operator version alone.}
    \item Furthermore, \textbf{$\mathbf{Int}_{\mathrm{op},\mathrm{ae}}$ sees
      neither the base space nor the realization as a function space}
      (Proposition~\ref{prop:Q-full-faithful},
      Corollary~\ref{cor:op-ae-iso-criterion}).
      Objects are distinguished only up to the pre-AL space
      $(\widetilde{L},\widetilde{I})$, and what the Daniell extension gives is
      its Banach completion
      (Proposition~\ref{prop:D-op-is-banach-completion}).
  \end{itemize}

  The last two points are the reasons for keeping
  Section~\ref{chap:functoriality} as an independent section.
  The first is a reason at the level of theorems and the second at the level of
  categories, and the latter is the stronger of the two.
  Only the side of Section~\ref{chap:functoriality} retains the distinction
  between an abstract completion and a concrete realization
  (Subsection~\ref{sec:prior-work}) as categorical structure.
  For details see Subsection~\ref{sec:what-op-ae-sees}.
\end{remark}
